\section{Three Connected Research Questions}
Before a collective AI decision, one agent's research can benefit both agents, although the researcher alone pays its cost. This creates a free-riding problem for AI-system designers and people relying on collective recommendations. Persico's committee model likewise studies costly information gathering before a group decision, with information as a public good; its voting environment differs from ours~\cite{persico2004}. Figure~\ref{fig:teaser} separates the completed formal benchmark from a future behavioral test.

\noindent\textbf{Q1---Economics:} How do private research costs shape strategic information acquisition and free-riding before a collective AI decision?
\textbf{Q2---Computation:} Can reproducible game-theoretic computation recover the predicted equilibria and show how equilibrium behavior changes as research costs vary?
\textbf{Q3---Behavior:} Do decision-makers follow equilibrium free-riding incentives, or can bounded reasoning and beliefs about others' willingness to research produce departures?
\textbf{Integrated:} When information is costly but collectively valuable, when will AI agents research rather than free-ride, and how closely will observed behavior match the formal prediction?
Economics supplies the incentive benchmark; computation makes it reproducible; behavior tests whether actual choices conform. The final comparison remains unanswered.
